\begin{ouverture}[calcul]{1}{Opérations et priorités}{Partie 1 · Nombres et calculs}
{Décimaux · Parenthèses · Division · Calcul réfléchi}
\objectif{Choisir une opération et en estimer le résultat.}
\objectif{Diviser par un nombre décimal non nul.}
\objectif{Traduire un problème par une expression.}
\prerequis{Tables de multiplication et numération décimale.}
\end{ouverture}
\pagestyle{otMain}\essentiel\fontsize{11}{15}\selectfont
\phantomsection\label{cours-1-1}
\essSection{1}{Organiser les calculs}
\begin{definition}Une somme réunit des termes ; un produit réunit des facteurs. Dans $7+3\times4$, la dernière opération est une addition : l’expression est une somme. Dans $(7+3)\times4$, c’est un produit.\end{definition}
\begin{propriete}Calculer d’abord les parenthèses, puis les multiplications et divisions, enfin les additions et soustractions. À priorité égale, calculer de gauche à droite.\end{propriete}
\begin{exemple}$18-6:3+2\times4=18-2+8=24$. Mais $(18-6):(3+2)\times4=12:5\times4=9{,}6$. Réécrire les parties inchangées à chaque étape.\end{exemple}
\begin{savoirfaire}[Anticiper et contrôler]Avant de calculer, prévoir un ordre de grandeur. Pour $19{,}8\times5{,}1$, $20\times5=100$ est un contrôle utile. Une estimation ne remplace pas le résultat exact.\end{savoirfaire}
\clearpage
\phantomsection\label{cours-1-2}
\essSection{2}{Diviser et calculer efficacement}
\begin{definition}Pour $b\ne0$, $a:b$ est le nombre qui, multiplié par $b$, donne $a$. Multiplier le dividende et le diviseur par un même nombre non nul conserve le quotient.\end{definition}
\begin{exemple}$7{,}56:0{,}6=75{,}6:6=12{,}6$. On transforme le diviseur en entier. Contrôle : $12{,}6\times0{,}6=7{,}56$. Diviser par $0{,}6$, qui est inférieur à 1, peut augmenter le nombre.\end{exemple}
\begin{propriete}$k(a+b)=ka+kb$ et $k(a-b)=ka-kb$. La distributivité permet de développer ou de regrouper un calcul : $17\times99=17(100-1)=1700-17=1683$.\end{propriete}
\begin{attention}Une division par zéro est impossible. Dans un problème, distinguer le nombre exact de parts et le nombre entier d’objets à acheter : $23:4=5{,}75$, mais il faut six boîtes de quatre places pour 23 objets.\end{attention}

\clearpage\rubrique{Découverte}
\begin{blocLivre}{ex-1-1}\begin{exercice}[1]{Addition décimale}
Calculer $4{,}7+2{,}85$.
\end{exercice}\end{blocLivre}
\begin{blocLivre}{ex-1-2}\begin{exercice}[1]{Différence décimale}
Calculer $10-3{,}48$.
\end{exercice}\end{blocLivre}
\begin{blocLivre}{ex-1-3}\begin{exercice}[1]{Produit de décimaux}
Calculer $0{,}8\times1{,}5$.
\end{exercice}\end{blocLivre}
\begin{blocLivre}{ex-1-4}\begin{exercice}[1]{Diviser par un décimal}
Calculer $9{,}6:0{,}4$.
\end{exercice}\end{blocLivre}
\begin{blocLivre}{ex-1-5}\begin{exercice}[1]{Nommer une expression}
Dans $5+2\times7$ puis $(5+2)\times7$, quelle est la dernière opération ? Calculer.
\end{exercice}\end{blocLivre}

\rubrique{Entraînement}
\begin{blocLivre}{ex-1-6}\begin{exercice}[2]{Même priorité}
Calculer $36:6\times3$ et $24-8+5$.
\end{exercice}\end{blocLivre}
\begin{blocLivre}{ex-1-7}\begin{exercice}[2]{Deux étages de parenthèses}
Calculer $42-[3\times(8-5)+7]$.
\end{exercice}\end{blocLivre}
\begin{blocLivre}{ex-1-8}\begin{exercice}[2]{Distribuer pour gagner du temps}
Calculer $23\times101$ et $48\times9{,}5+48\times0{,}5$.
\end{exercice}\end{blocLivre}
\begin{blocLivre}{ex-1-9}\begin{exercice}[2]{Une division à contrôler}
Calculer $6{,}3:0{,}15$ puis contrôler le résultat.
\end{exercice}\end{blocLivre}
\begin{blocLivre}{ex-1-10}\begin{exercice}[2]{Corriger une priorité}
Léo écrit $7+3\times5=10\times5=50$. Expliquer son erreur.
\end{exercice}\end{blocLivre}

\rubrique{Problèmes}
\begin{blocLivre}{ex-1-11}\begin{exercice}[3]{Une sortie}
Quatre élèves paient chacun 6,50 euros de transport et 3,20 euros d’entrée. Écrire une seule expression pour le total puis calculer.
\end{exercice}\end{blocLivre}
\begin{blocLivre}{ex-1-12}\begin{exercice}[3]{Des bouteilles}
On verse 8,4 L dans des bouteilles de 0,75 L. Combien remplit-on de bouteilles entières ? Quel volume reste-t-il ?
\end{exercice}\end{blocLivre}
\begin{blocLivre}{ex-1-13}\begin{exercice}[3]{Un budget partagé}
Une association a 125 euros. Elle achète cinq lots à 12,40 euros et partage le reste également entre six ateliers. Quelle somme reçoit chaque atelier ?
\end{exercice}\end{blocLivre}
\begin{blocLivre}{ex-1-14}\begin{exercice}[3]{Une expression à inventer}
Avec 8, 6 et 2 utilisés chacun une fois, proposer un calcul donnant 28. Justifier le rôle des parenthèses.
\end{exercice}\end{blocLivre}
\begin{blocLivre}{ex-1-15}\begin{exercice}[3]{Comparer deux offres}
Un magasin vend un carnet 2,80 euros. Un autre vend un lot de cinq carnets 12 euros avec 1,50 euro de frais pour toute commande. Pour dix carnets, comparer les coûts.
\end{exercice}\end{blocLivre}

\clearpage
